% Generated draft: requires mathematical and source review.
\documentclass[11pt]{article}
\usepackage[margin=25mm]{geometry}
\usepackage{fontspec}
\setmainfont{Latin Modern Roman}
\usepackage{amsmath,amssymb,mathtools,unicode-math}
\setmathfont{Latin Modern Math}
\usepackage{enumitem,needspace,xurl,hyperref}
\setlength{\parindent}{0pt}
\setlength{\parskip}{5pt}
\setlength{\emergencystretch}{3em}
\newcommand{\N}{\mathbb{N}}
\newcommand{\Z}{\mathbb{Z}}
\newcommand{\Q}{\mathbb{Q}}
\newcommand{\R}{\mathbb{R}}
\newcommand{\C}{\mathbb{C}}
\newcommand{\F}{\mathbb{F}}
\newcommand{\A}{\mathbb{A}}
\newcommand{\PP}{\mathbb{P}}
\newcommand{\E}{\mathbb{E}}
\newcommand{\eps}{\varepsilon}
\newcommand{\dd}{\,\mathrm{d}}
\begin{document}
\section*{UnsolvedMath ID: 30006446}
\textbf{Source problem code:} catalogue entry
\subsection{Definitions and mathematical statement}
Retain the source notation \(A=k[x_1,x_2]\); the supplied record does not specify hypotheses on the coefficient object \(k\). A monomial ideal \(I\subseteq A\) is an ideal generated by monomials \(x_1^a x_2^b\), where \(a,b\in\Z_{\geq 0}\).

\emph{Conjectural question supplied by the record.} For every realizable monomial ideal \(I\subseteq k[x_1,x_2]\), is the homological-dimension lower bound for the number of realizations of \(I\) sharp?

\emph{Unresolved formulation.} The supplied excerpt begins with the assertion that ``this bound'' is conjecturally sharp, but omits the preceding bound and its definitions. In particular, it does not define realizability, the objects counted as realizations, their equivalence relation if any, the homological dimension involved, or the formula for the lower bound. Thus the question above preserves the supplied universal quantifier over realizable monomial ideals, but cannot be converted into a self-contained numerical equality from this record. The hypotheses on \(k\), any restrictions on \(I\), and the precise meaning of sharpness require recovery from the omitted context.
\subsection{Short English statement}
The supplied conjecture asks whether a lower bound expressed using homological dimension is sharp for every realizable monomial ideal in \(k[x_1,x_2]\). The record omits the bound and the definitions of realization and realizability, so a precise numerical formulation remains unresolved.
\subsection{Sources}
\begin{itemize}
\item Supplied local UnsolvedMath registry record: permanent ID \(30006446\), title \emph{Sharp Homological Bounds for Realizable Monomial Ideals}, exact source code \texttt{catalogue entry}, in the Oberwolfach Reports Open Problems collection. No proposer is named; the supplied proposed year is 2025.
\item Supplied source citation: \emph{Singularities (2026)}, Oberwolfach Reports, pp.\ 2413--2490, DOI \texttt{10.4171/owr/2025/45}. Supplied source and literature link: \url{https://doi.org/10.4171/owr/2025/45}, also labelled \emph{Oberwolfach Report 45/2025}. The relationship between the title year 2026 and the report designation 2025 is not explained in the record.
\item Literature assessment dated exactly \texttt{2026-08-22}: the supplied curation labels the problem open and states that the homological-dimension lower bound is proved, while its sharpness for realizable monomial ideals in two variables is explicitly conjectural. This is supplied dated curation, not an independently verified literature assessment.
\item The registry labels its statement \texttt{corrected\_verified} and reports that the official report was checked and the question isolated with repaired notation. These are supplied verification claims; no independent source inspection is asserted here.
\item The surrounding excerpt attributes analytic lattice homology of normal surface singularities and reduced curve singularities to \mbox{\'Agoston} and N\'emethi through references [1] and [2]. Their bibliographic entries are not supplied, and this contextual attribution does not identify the proposer of the conjecture.
\end{itemize}
\end{document}
